\section{Cumulus parameterization}
{\bf \Large
\begin{tabular}{ccc}
\hline
  Corresponding author & : & Yuta Kawai\\
\hline
\end{tabular}
}

Cumulus parameterization represents the effects of unresolved moist convection on the resolved atmospheric state.
In SCALE-DG, a moist convective adjustment scheme is available as a simple representation of cumulus convection.

\subsection{Moist convective adjustment}
\label{subsec:moist_convective_adjustment}
%%
The scheme is based on the concept of \citet{manabe1965}.
Saturated convective instability is diagnosed from the vertical distribution of saturated moist static energy (MSE) in nearly saturated layers.
If such instability is detected, the unstable layer is adjusted toward a moist-neutral thermodynamic profile
subject to column-integrated MSE conservation and the available-water constraint.
Water vapor condensed during the adjustment is immediately removed as precipitation.
The adjustment is repeated until no saturated convective instability remains in the column.
The overall procedure of the adjustment is summarized in Fig.~\ref{fig:mca_procedure}.


The saturated MSE is defined as
%
\begin{align}
 h_s = C_{p}(q_{v,s}) T + gz + L q_{v,s},
\label{eq:mca_hsat}
\end{align}
%
where
$C_{p}(q_v) = (1-q_v) C_{p,d} + q_v C_{p,v}$
and $q_{v,s}$ is the saturation specific humidity.
A layer is regarded as potentially convectively unstable when it is nearly saturated and $h_s$ decreases with height, i.e.,
%
\begin{align}
 \frac{\partial h_s}{\partial z} < 0.
 \label{eq:mca_instability}
\end{align}
%
In practice, small thresholds are introduced for relative humidity and the vertical difference of $h_s$
in order to avoid triggering the adjustment by numerical noise.
The contiguous layer satisfying these criteria is referred to as the \textit{diagnosed unstable core} $\mathcal{C}$.
The \textit{adjustment region} $\mathcal{A}$ is the layer over which the moist-neutral profile and the column-integrated constraints are imposed.
It initially corresponds to $\mathcal{C}$ but may be extended vertically when necessary to obtain a water-feasible and energy-conserving adjusted state.

For a given $\mathcal{A}$, the adjusted thermodynamic profile is determined using the temperature $T_b$ at the base of $\mathcal{A}$ as a single scalar degree of freedom.
For each trial value of $T_b$, a moist-neutral temperature profile is constructed by requiring vertically uniform saturated MSE:
%
\begin{align}
 h_s(\widetilde{T}(z;T_b),\widetilde{p}(z;T_b),z) = h_{s,0},
 \label{eq:mca_neutral}
\end{align}
%
where a tilde denotes a trial state before the diagnosed precipitation is removed.
Equation~\eqref{eq:mca_neutral} is integrated upward from $T_b$ to obtain $\widetilde{T}(z;T_b)$.

The corresponding water-vapor profile is constructed subject to the saturation constraint.
Nodes that are nearly saturated within $\mathcal{C}$ are adjusted to saturation,
%
\begin{align}
 \widetilde{q}_v(z;T_b)
 =
 q_{v,s}\left(\widetilde{T}(z;T_b),\widetilde{p}(z;T_b)\right),
\end{align}
%
whereas at the other nodes within $\mathcal{A}$,
%
\begin{align}
 \widetilde{q}_v(z;T_b)
 =
 \min \left[
 q_v^*,
 q_{v,s}\left(\widetilde{T}(z;T_b),\widetilde{p}(z;T_b)\right)
 \right].
\label{eq:mca_qv}
\end{align}
%
Here, the asterisk denotes the state before the adjustment.
During this trial-profile construction, the density $\rho^*$ is held fixed, while the pressure is updated through
%
\begin{align*}
 \widetilde{p} = \rho^* R(\widetilde{q}_v)\widetilde{T}, \qquad
 R(q_v) = (1-q_v)R_d + q_v R_v.
\end{align*}
%
The moist-neutral profile and pressure are determined consistently by iteration.

For each $T_b$, the column-integrated water vapor and MSE of the resulting trial profile are
%
\begin{align*}
 W(T_b)
 =
 \int_{\mathcal{A}} \rho^* \widetilde{q}_v(z;T_b)\,dz ,
 \qquad
 H(T_b)
 =
 \int_{\mathcal{A}} \rho^*
 \left[
 C_p(\widetilde{q}_v)\widetilde{T}
 + gz
 + L\widetilde{q}_v
 \right] dz ,
\end{align*}
%
and their values before the adjustment are
%
\begin{align*}
 W^*
 =
 \int_{\mathcal{A}} \rho^* q_v^*\,dz ,
 \qquad
 H^*
 =
 \int_{\mathcal{A}} \rho^*
 \left[
 C_p(q_v^*)T^*
 + gz
 + Lq_v^*
 \right] dz .
\end{align*}
%
The available-water constraint first restricts the admissible range of $T_b$ according to
%
\begin{align}
 W(T_b) \le W^*.
 \label{eq:mca_water_constraint}
\end{align}
%
Within this range, $T_b$ is determined by imposing conservation of the column-integrated MSE:
%
\begin{align}
 H(T_b) = H^*.
 \label{eq:mca_energy_constraint}
\end{align}
%
Thus, the water constraint restricts the admissible family of trial profiles, and the MSE constraint selects the final value of $T_b$.
If no such solution exists within the initially diagnosed unstable core, $\mathcal{A}$ is progressively extended vertically and the calculation is repeated.
This extension is a numerical procedure for obtaining a consistent adjusted state and does not represent physical entrainment.

Once $T_b$ has been determined, the reduction in column-integrated water vapor is immediately removed as precipitation:
%
\begin{align}
 P_{\rm MCA}
 =
 W^* - W(T_b)
 =
 \int_{\mathcal{A}}
 \rho^*
 \left(q_v^*-\widetilde{q}_v\right) dz
 \ge 0.
 \label{eq:mca_precip}
\end{align}
%
The precipitation mass is distributed within $\mathcal{A}$ over nodes where the local water-vapor mass density decreases.
No condensate is retained as a prognostic cloud-water variable.


After the state has been updated, saturated convective instability is diagnosed again over the entire column.
If another unstable core is found, the adjustment procedure described above is repeated using the updated state.
This diagnosis--adjustment cycle continues until no saturated convective instability remains in the column.
The precipitation produced in successive adjustments is accumulated throughout this cycle.


Let $\rho_{\rm precip}(z)$ denote the accumulated local precipitation mass density after completion of the diagnosis--adjustment cycle.
The final density is obtained as
%%
\begin{align*}
\rho^{\rm adj} = \rho^* - \rho_{\rm precip}.
\end{align*}
%%
The specific humidity changes consistently with the reduced density, 
while the water-vapor mass density remains unchanged during this final precipitation-removal step.
The pressure, potential temperature, and other thermodynamic quantities are then recomputed consistently to obtain the final adjusted state.


The adjustment is applied over one physics time step $\Delta t_{\rm MCA}$.
The tendencies are
%
\begin{align}
 S_{\rho,{\rm MCA}}
 &=
 \frac{\rho^{\rm adj}-\rho^*}{\Delta t_{\rm MCA}},
 \\
 S_{\rho\theta,{\rm MCA}}
 &=
 \frac{ \rho^{\rm adj}\theta^{\rm adj} - \rho^*\theta^* } {\Delta t_{\rm MCA}},
 \\
 S_{\rho q_v,{\rm MCA}}
 &=
 \frac{ \rho^{\rm adj}q_v^{\rm adj} - \rho^*q_v^* } {\Delta t_{\rm MCA}} .
 \label{eq:mca_tendencies}
\end{align}
%
The precipitation mass flux and the associated energy flux removed from the atmospheric column are diagnosed from the precipitation due to the adjustment.


%%%%%%%%%%%%%%%%%%%%%%%%%%%
\begin{figure}[htbp]
\centering
\begin{tikzpicture}[
  font=\scriptsize,
  scale=0.8,
  node distance=3mm,
  box/.style={
    rectangle,
    rounded corners,
    draw=black,
    align=center,
    text width=5.6cm,
    minimum height=6mm,
    inner sep=2.5pt
  },
  decision/.style={
    diamond,
    draw=black,
    align=center,
    aspect=3.2,
    inner sep=0.5pt
  },
  arrow/.style={
    ->,
    >=Stealth,
    semithick
  }
]

\node[box] (initial)
{Initial atmospheric state};

\node[box, below=of initial] (diagnose)
{Diagnose saturated convective instability\\
and determine the unstable core $\mathcal{C}$};

\node[decision, below=of diagnose] (unstable)
{Unstable core exists?};

\node[box, below=5mm of unstable] (region)
{Define the adjustment region $\mathcal{A}$};

\node[box, below=of region] (trial)
{For a trial $T_b$, construct the moist-neutral profile\\
$\widetilde{T}(z;T_b)$, $\widetilde{q}_v(z;T_b)$, and
$\widetilde{p}(z;T_b)$};

\node[box, below=of trial] (water)
{Restrict the admissible range of $T_b$\\
by $W(T_b)\le W^*$};

\node[decision, below=4mm of water] (waterfeasible)
{Water-feasible range exists?};

\node[box, below=5mm of waterfeasible] (energy)
{Determine $T_b$ within the feasible range\\
from $H(T_b)=H^*$};

\node[decision, below=4mm of energy] (energysol)
{Energy-conserving solution exists?};

\node[box, right=20mm of energysol, text width=3.0cm] (extend)
{Extend $\mathcal{A}$ vertically};

\node[box, below=5mm of energysol] (precip)
{Diagnose and accumulate precipitation};

\node[box, below=of precip] (update)
{Update the atmospheric state};

\node[box, right=22mm of unstable, text width=4.0cm] (final)
{Construct the final adjusted state\\
using accumulated precipitation};

\node[box, below=of final, text width=4.0cm] (tendency)
{Evaluate MCA tendencies over $\Delta t_{\rm MCA}$};

\draw[arrow] (initial) -- (diagnose);
\draw[arrow] (diagnose) -- (unstable);

\draw[arrow]
  (unstable)
  -- node[right, font=\scriptsize] {yes}
  (region);

\draw[arrow]
  (unstable.east)
  -- node[above, font=\scriptsize] {no}
  (final.west);

\draw[arrow] (region) -- (trial);
\draw[arrow] (trial) -- (water);
\draw[arrow] (water) -- (waterfeasible);

\draw[arrow]
  (waterfeasible)
  -- node[right, font=\scriptsize] {yes}
  (energy);

\draw[arrow] (energy) -- (energysol);

\draw[arrow]
  (energysol)
  -- node[right, font=\scriptsize] {yes}
  (precip);

\draw[arrow]
  (waterfeasible.east)
  -- node[above, font=\scriptsize] {no}
  ++(11mm,0)
  |- (extend.north);

\draw[arrow]
  (energysol.east)
  -- node[above, font=\scriptsize] {no}
  (extend.west);

\draw[arrow]
  (extend.north)
  |- (trial.east);

\draw[arrow] (precip) -- (update);

\draw[arrow]
  (update.west)
  -- ++(-10mm,0)
  |- (diagnose.west);

\draw[arrow] (final) -- (tendency);

\end{tikzpicture}
\caption{
Schematic procedure of the moist convective adjustment scheme.
The available-water constraint first restricts the admissible range of $T_b$,
and the MSE constraint then determines $T_b$ within this range.
The diagnosis--adjustment cycle is repeated until no saturated convective
instability remains in the column.
}
\label{fig:mca_procedure}
\end{figure}